\documentclass[addpoints,spanish, 12pt,a4paper]{exam}
%\documentclass[answers, spanish, 12pt,a4paper]{exam}
% \printanswers
\renewcommand*\half{.5}

\pointpoints{punto}{puntos}
\hpword{Puntos:}
\vpword{Puntos:}
\htword{Total}
\vtword{Total}
\hsword{Resultado:}
\hqword{Ejercicio:}
\vqword{Ejercicio:}

\usepackage{booktabs} 
\usepackage[utf8]{inputenc}
\usepackage[spanish]{babel}
\usepackage{eurosym}
%\usepackage[spanish,es-lcroman, es-tabla, es-noshorthands]{babel}


\usepackage[margin=1in]{geometry}
\usepackage{amsmath,amssymb}
\usepackage{multicol}
\usepackage{yhmath}

\pointsinrightmargin % Para poner las puntuaciones a la derecha. Se puede cambiar. Si se comenta, sale a la izquierda.
\extrawidth{-2.4cm} %Un poquito más de margen por si ponemos textos largos.
\marginpointname{ \emph{\points}}

\usepackage{graphicx}

\graphicspath{{../img/}} 

\newcommand{\class}{1 ESO}
\newcommand{\examdate}{\today}
\newcommand{\examnum}{Álgebra}
\newcommand{\tipo}{A}


\newcommand{\timelimit}{45 minutos}

\renewcommand{\solutiontitle}{\noindent\textbf{Solución:}\enspace}


\pagestyle{head}
\firstpageheader{\includegraphics[width=0.2\columnwidth]{header_left}}{\textbf{Departamento de Matemáticas\linebreak \class}\linebreak \examnum}{\includegraphics[width=0.1\columnwidth]{header_right}}
\runningheader{\class}{\examnum}{Página \thepage\ de \numpages}
\runningheadrule


\usepackage{pgf,tikz,pgfplots}
\pgfplotsset{compat=1.15}
\usepackage{mathrsfs}
\usetikzlibrary{arrows}


\begin{document}

\noindent
\begin{tabular*}{\textwidth}{l @{\extracolsep{\fill}} r @{\extracolsep{6pt}} }
\textbf{Nombre:} \makebox[3.5in]{\hrulefill} & \textbf{Fecha:}\makebox[1in]{\hrulefill} \\
 & \\
\textbf{Tiempo: \timelimit} & Tipo: \tipo 
\end{tabular*}
\rule[2ex]{\textwidth}{2pt}
\textbf{Instrucciones:} Justifica todos los
resultados.
Esta prueba tiene \numquestions\ ejercicios con una puntuación máxima de \numpoints. 
La nota del examen se calculará de manera proporcional a la puntuación obtenida. 
% \textbf{Para recuperar el examen anterior se tendrán en cuenta las preguntas 1 a 3}\\


\begin{center}


\addpoints
 %\gradetable[h][questions]
	\pointtable[h][questions]
   
\end{center}
% \textbf{Nota:} Se tendrán en cuenta los ejercicios 1 y 2 para recuperar contenidos previos.

\noindent
\rule[2ex]{\textwidth}{2pt}



\begin{questions}

% \question
% Expresa de forma algebraica los siguientes enunciados matemáticos:
% \begin{parts}
% \part[1] La suma de un número, \( a \), y su mitad.
% \vspace{30pt}
% \part[1] El triple de la mitad de un número, \( n \).
% \vspace{30pt}
% \part[1] El área de un cuadrado de lado \( a \).
% \vspace{30pt}
% \end{parts}

% \question[1]
% Rodea con un círculo aquellas expresiones algebraicas que sean monomios:
% \[
% 6x^3 + 3y^4 \quad\quad 6ab \quad\quad 5xyz \quad\quad 7y^5 + 4x^3 \quad\quad 2y^3
% \]
% \vspace{30pt}

\question[2]
Completa la tabla indicando el coeficiente, la parte literal y el grado de cada monomio:
\vspace{10pt}

\begin{tabular}{|c|c|c|c|}
\hline
\textbf{Monomio} & \textbf{Coeficiente} & \textbf{Parte literal} & \textbf{Grado} \\
\hline
3 & & & \\
\hline
$9x^2$ & & & \\
\hline
$\frac{2}{3}bc$ & & & \\
\hline
$-ab^3x$ & & & \\
\hline
\end{tabular}


\question[1]
Rodea con un círculo los monomios que sean semejantes:
\[
3a^2b^3x \quad\quad 4a^2b^3x \quad\quad -2a^2b^3x \quad\quad 5a^2b^2x^2
\]

\question[1]
Realiza las siguientes operaciones con monomios:
\begin{multicols}{2}
\begin{parts}
\part \( 3x^2 + 5x^2 \)
\vspace{30pt}
\part \( 7a^3 - 2a^3 \)
\vspace{30pt}
\part \( (2x)(-4x^3) \)
\vspace{30pt}
\part $\frac{6x^5}{2x^2}$
\vspace{30pt}
\end{parts}
\end{multicols}

\question[1]
Calcula el valor numérico de la siguiente expresión algebraica sabiendo que \( x = 2 \) y \( y = -3 \):
\[
4x^2 - 2xy + y^2
\]
\vspace{30pt}

\question
Rodea, en cada caso, el valor de \( x \) que es solución de la ecuación:
\begin{parts}
\part[1] \( 2x + 5 = 25 \quad \Rightarrow \quad x = 5 \quad x = 10 \quad x = 15 \quad x = 20 \)
\part[1] \( 3x - 4 = 14 \quad \Rightarrow \quad x = 2 \quad x = 4 \quad x = 6 \quad x = 8 \)
\end{parts}

\question
Resuelve las siguientes ecuaciones:
\begin{parts}
\part[1] \( x + 6 = 15 \)
\vspace{30pt}
\part[1] \( x - 4 = 9 \)
\vspace{30pt}
\part[1] \( \frac{x}{3} = 2 \)
\vspace{30pt}
\part[1] \( \frac{x}{2} = 12 \)
\vspace{30pt}
\end{parts}

\question
Resuelve las siguientes ecuaciones:
\begin{parts}
\part[1] \( x + 13 = 2x + 11 \)
\vspace{30pt}
\part[1] \( x + 10 = 3x + 4 \)
\vspace{30pt}
\end{parts}

\question
Resuelve las siguientes ecuaciones:
\begin{parts}
\part[1] \( (5x - 3) + (7x + 6) = 13 \)
\vspace{80pt}
\part[1] \( (2x + 4) + (1 - x) = 14 \)
\vspace{80pt}
\end{parts}

% \question[1]
% La suma de tres números consecutivos es 42. ¿Cuáles son esos números?
% \vspace{30pt}

% \question[1]
% Juan tiene 25 euros más que Mario y 30 euros menos que Enrique. ¿Cuánto tiene cada uno sabiendo que entre los tres tienen 140 euros?
% \vspace{30pt}





\end{questions}



\end{document}
\grid
